

\begin{derivation}[从\eqnrefs{eq:heat-kernel-a4-flat-dp11}到\eqnrefs{eq:qcd-glue-ghost-uv-coefficient-dp10}]
纯 Yang--Mills 一圈系数中的 $11/3$ 来自规范矢量涨落与 Faddeev--Popov 鬼场对同一个背景场强平方项的不同热核贡献。其定量来源由规范固定后的二次算符分解给出。把完整规范联络分成经典背景场与量子涨落，

\begin{equation*}
 \mathscr C_\mu=\bar{\mathscr C}_\mu+Q_\mu.
\end{equation*}
背景 Feynman 规范取
\begin{equation*}
 \mathscr C_{\rm bg}^a[Q]=(\bar D^\mu Q_\mu)^a,
 \qquad\xi=1.
\end{equation*}
把 Yang--Mills 作用量与规范固定项展开到 $Q^2$；背景场满足经典运动方程时线性项消失，二次项可写为
\begin{equation*}
 S_E^{(2)}
 =\frac\hbar2\int\dd^4 x\,
 Q_\mu^a(\Delta_1)_{\mu\nu}^{ab}Q_\nu^b
 +\hbar\int\dd^4 x\,
 \bar c^a(\Delta_0)^{ab}c^b,
\end{equation*}
其中欧几里得算符为
\begin{align*}
 (\Delta_1)_{\mu\nu}
 &=-\bar D^2\delta_{\mu\nu}
 -2\ii g_s\bar{\mathcal F}_{\mu\nu}^aT_{\rm adj}^a,\\
 \Delta_0&=-\bar D^2.
\end{align*}
这里 $(T_{\rm adj}^a)_{bc}=-\ii f^{abc}$，所以
\begin{equation*}
 \Tr_{\rm adj}(T_{\rm adj}^aT_{\rm adj}^b)
 =C_A\delta^{ab}.
\end{equation*}
高斯积分给出纯规范扇区的一圈有效作用量
\begin{equation*}
 \Gamma_{g+\rm gh}^{(1)}
 =\frac12\Tr\ln\Delta_1-\Tr\ln\Delta_0.
\end{equation*}
第二项没有 $1/2$，因为 $c,\bar c$ 是彼此独立的复 Grassmann 场。

把这些算符逐项代入 \eqnrefs{eq:heat-kernel-a4-flat-dp11}。对背景伴随表示协变导数，丛曲率为
\begin{equation*}
 \Omega_{\rho\sigma}
 =[\bar D_\rho,\bar D_\sigma]
 =\ii g_s\bar{\mathcal F}_{\rho\sigma}^aT_{\rm adj}^a.
\end{equation*}
先做颜色迹：
\begin{align*}
 \Tr_{\rm adj}(\Omega_{\rho\sigma}\Omega_{\rho\sigma})
 &=(\ii g_s)^2
 \bar{\mathcal F}_{\rho\sigma}^a\bar{\mathcal F}_{\rho\sigma}^b
 \Tr_{\rm adj}(T^aT^b),\\
 &=-g_s^2C_A\bar{\mathcal F}_{\rho\sigma}^a
 \bar{\mathcal F}_{\rho\sigma}^a.
\end{align*}
矢量算符还需对四个 Lorentz 分量取迹，因此
\begin{equation*}
 \tr_{V\otimes\rm adj}\Omega^2
 =-4g_s^2C_A\bar{\mathcal F}^2.
\end{equation*}

与 Laplace 型标准形式 $\Delta=-(\nabla^2+E)$ 比较得到
\begin{equation*}
 (E_1)_{\mu\nu}
 =2\ii g_s\bar{\mathcal F}_{\mu\nu}^aT_{\rm adj}^a.
\end{equation*}
计算 $E_1^2$ 时，Lorentz 反对称性给
\begin{equation*}
 \bar{\mathcal F}_{\mu\nu}\bar{\mathcal F}_{\nu\mu}
 =-\bar{\mathcal F}_{\mu\nu}\bar{\mathcal F}_{\mu\nu}.
\end{equation*}
所以两个负号相消：
\begin{align*}
 \tr E_1^2
 &=(2\ii g_s)^2
 \bar{\mathcal F}_{\mu\nu}^a\bar{\mathcal F}_{\nu\mu}^b
 \Tr_{\rm adj}(T^aT^b),\\
 &=4g_s^2C_A\bar{\mathcal F}^2.
\end{align*}
因此矢量热核的对数局域系数为
\begin{align*}
 a_4(\Delta_1)
 &\supset
 \frac12\tr E_1^2+\frac1{12}\tr\Omega^2,\\
 &=\left(2-\frac13\right)g_s^2C_A\bar{\mathcal F}^2,\\
 &=\frac53g_s^2C_A\bar{\mathcal F}^2.
\end{align*}
对鬼场，$E_0=0$ 且没有矢量分量的重数：
\begin{equation*}
 a_4(\Delta_0)
 \supset-\frac1{12}g_s^2C_A\bar{\mathcal F}^2.
\end{equation*}
再乘上行列式对应的权重：
\begin{align*}
 \frac12a_4(\Delta_1)-a_4(\Delta_0)
 &=\frac12\frac53g_s^2C_A\bar{\mathcal F}^2
 -\left(-\frac1{12}g_s^2C_A\bar{\mathcal F}^2\right),\\
 &=\frac{11}{12}g_s^2C_A\bar{\mathcal F}^2.
\end{align*}
经典动能算符对应的拉格朗日量项为 $\bar{\mathcal F}^2/4$，因此把极点系数改写成相对于 $\frac14\bar{\mathcal F}^2$ 的系数需要乘以 $4$：
\begin{equation*}
 {4\times\frac{11}{12}C_A=\frac{11}{3}C_A.}
\end{equation*}
这就是 \eqnrefs{eq:qcd-glue-ghost-uv-coefficient-dp10} 的 $11C_A/3$。它不是某一张胶子圈图的独立数值，而是矢量涨落与 Faddeev--Popov 行列式在同一背景规范不变算符上的总和。
\end{derivation}

\begin{derivation}[从\eqnrefs{eq:reduced-loop-measure-si-r10}到\eqnrefs{eq:qcd-quark-loop-uv-coefficient-dp10}]
取质量为 $m_f$、属于颜色表示 $R_f$ 的一个 Dirac 费米子。颜色生成元满足
\begin{equation*}
 \Tr_{R_f}(T^aT^b)=T(R_f)\delta^{ab}.
\end{equation*}
采用统一的约化物理四动量圈积分记号，闭合费米子环的二点函数写成
\begin{equation*}
 \ii\Pi_{f,ab}^{\mu\nu}(q)
 =-g_s^2T(R_f)\delta_{ab}
 \int_\ell^{(d)}
 \frac{N^{\mu\nu}(\ell,q)}
 {D_0D_1},
 \qquad d=4-2\epsilon,
\end{equation*}
其中
\begin{align*}
 D_0&=\ell^2-m_f^2c^2+\ii0,\\
 D_1&=(\ell+q)^2-m_f^2c^2+\ii0,\\
 N^{\mu\nu}(\ell,q)
 &=\Tr\!\left[
 \gamma^\mu(\slashed\ell+m_fc)
 \gamma^\nu(\slashed\ell+\slashed q+m_fc)
 \right].
\end{align*}

先证明横向性，而不是在积分后凭 Lorentz 协变性猜结构。定义去掉整体 $\ii$ 的自由 Dirac 传播核
\begin{equation*}
 \mathsf S(\ell)
 \equiv\frac{\slashed\ell+m_fc}{D_0}
 =(\slashed\ell-m_fc)^{-1},
\end{equation*}
以及
\begin{equation*}
 \mathsf S(\ell+q)
 =\frac{\slashed\ell+\slashed q+m_fc}{D_1}.
\end{equation*}
于是
\begin{equation*}
 \slashed q
 =\mathsf S^{-1}(\ell+q)-\mathsf S^{-1}(\ell).
\end{equation*}
把 $q_\mu\gamma^\mu=\slashed q$ 代入循环迹，并只在迹内部使用循环性：
\begin{align*}
 &\Tr\!\left[
 \slashed q\,\mathsf S(\ell)\gamma^\nu\mathsf S(\ell+q)
 \right]\\
 &\quad=
 \Tr\!\left[
 \mathsf S^{-1}(\ell+q)\mathsf S(\ell)
 \gamma^\nu\mathsf S(\ell+q)
 \right]
 -\Tr\!\left[
 \mathsf S^{-1}(\ell)\mathsf S(\ell)
 \gamma^\nu\mathsf S(\ell+q)
 \right]\\
 &\quad=
 \Tr\!\left[
 \mathsf S(\ell)\gamma^\nu
 \mathsf S(\ell+q)\mathsf S^{-1}(\ell+q)
 \right]
 -\Tr\!\left[
 \gamma^\nu\mathsf S(\ell+q)
 \right]\\
 &\quad=
 \Tr\!\left[\mathsf S(\ell)\gamma^\nu\right]
 -\Tr\!\left[\gamma^\nu\mathsf S(\ell+q)\right].
\end{align*}
因此由起始定义逐项收缩可得
\begin{align*}
 \ii q_\mu\Pi_{f,ab}^{\mu\nu}(q)
 &=-g_s^2T(R_f)\delta_{ab}
 \int_\ell^{(d)}
 \Tr\!\left[
 \slashed q\,\mathsf S(\ell)\gamma^\nu\mathsf S(\ell+q)
 \right]\\
 &=-g_s^2T(R_f)\delta_{ab}
 \int_\ell^{(d)}
 \left\{
 \Tr[\mathsf S(\ell)\gamma^\nu]
 -\Tr[\gamma^\nu\mathsf S(\ell+q)]
 \right\}.
\end{align*}
第二项作 $r=\ell+q$ 平移。维数正规化保持积分测度平移不变，且 $\Tr[\mathsf S(r)\gamma^\nu]=\Tr[\gamma^\nu\mathsf S(r)]$，所以两项积分完全相消：
\begin{equation*}
 q_\mu\Pi_{f,ab}^{\mu\nu}(q)=0.
\end{equation*}
颜色对角性与 Lorentz 协变性先给出 $\Pi_{f,ab}^{\mu\nu}=\delta_{ab}[A_f(q^2)\eta^{\mu\nu}+B_f(q^2)q^\mu q^\nu]$；再与上式联立得到 $A_f=-q^2B_f$，因此存在唯一标量函数 $\Pi_f(q^2)$ 使
\begin{equation*}
 \Pi_{f,ab}^{\mu\nu}(q)
 =\delta_{ab}(q^2\eta^{\mu\nu}-q^\mu q^\nu)\Pi_f(q^2).
\end{equation*}

接着真正计算系数。Dirac 迹展开为
\begin{align*}
 N^{\mu\nu}
 &=\Tr(\gamma^\mu\slashed\ell\gamma^\nu\slashed\ell)
 +\Tr(\gamma^\mu\slashed\ell\gamma^\nu\slashed q)
 +m_f^2c^2\Tr(\gamma^\mu\gamma^\nu)\\
 &=4\left[
 \ell^\mu(\ell+q)^\nu
 +\ell^\nu(\ell+q)^\mu
 -\eta^{\mu\nu}\ell\cdot(\ell+q)
 +m_f^2c^2\eta^{\mu\nu}
 \right].
\end{align*}
奇数个 $\gamma$ 矩阵的两项为零。用 Feynman 参数
\begin{equation*}
 \frac1{D_0D_1}
 =\int_0^1\dd x\,
 \frac1{[xD_0+(1-x)D_1]^2}
\end{equation*}
并定义
\begin{equation*}
 r\equiv\ell+(1-x)q,
 \qquad
 \Delta_f(x)\equiv m_f^2c^2-x(1-x)q^2.
\end{equation*}
则
\begin{equation*}
 xD_0+(1-x)D_1=r^2-\Delta_f(x)+\ii0.
\end{equation*}
把
\begin{equation*}
 \ell=r-(1-x)q,
 \qquad
 \ell+q=r+xq
\end{equation*}
代入分子，保留关于 $r$ 的偶次项：
\begin{align*}
 \frac14N^{\mu\nu}_{\rm even}
 ={}&2r^\mu r^\nu-\eta^{\mu\nu}r^2
 -2x(1-x)q^\mu q^\nu\\
 &+\eta^{\mu\nu}
 \left[m_f^2c^2+x(1-x)q^2\right].
\end{align*}
所有关于 $r$ 的奇次项在对称积分中为零。

关键的张量积分不直接用“对称性可得”跳过。维数正规化下全导数积分为零：
\begin{equation*}
 0=\int_r^{(d)}
 \frac{\partial}{\partial r_\mu}
 \left[
 \frac{r_\nu}{r^2-\Delta_f+\ii0}
 \right].
\end{equation*}
展开导数得到
\begin{equation*}
 \int_r^{(d)}
 \frac{2r^\mu r^\nu}
 {(r^2-\Delta_f+\ii0)^2}
 =\eta^{\mu\nu}
 \int_r^{(d)}
 \frac1{r^2-\Delta_f+\ii0}.
\end{equation*}
另一方面
\begin{align*}
 \frac{r^2}{(r^2-\Delta_f)^2}
 &=\frac1{r^2-\Delta_f}
 +\frac{\Delta_f}{(r^2-\Delta_f)^2}.
\end{align*}
因此 $r^\mu r^\nu$ 与 $r^2$ 两项组合成
\begin{align*}
 &\int_r^{(d)}
 \frac{2r^\mu r^\nu-\eta^{\mu\nu}r^2}
 {(r^2-\Delta_f+\ii0)^2}\\
 &\qquad=-\eta^{\mu\nu}\Delta_f
 \int_r^{(d)}
 \frac1{(r^2-\Delta_f+\ii0)^2}.
\end{align*}
再加上分子中的显式质量与 $q^2$ 项，利用
\begin{equation*}
 -\Delta_f+m_f^2c^2+x(1-x)q^2
 =2x(1-x)q^2,
\end{equation*}
整个偶分子积分化为唯一横向组合
\begin{align*}
 \int_r^{(d)}\frac{N^{\mu\nu}_{\rm even}}
 {(r^2-\Delta_f+\ii0)^2}
 =8x(1-x)
 (q^2\eta^{\mu\nu}-q^\mu q^\nu)
 \int_r^{(d)}\frac1{(r^2-\Delta_f+\ii0)^2}.
\end{align*}

剩下的标量积分也直接算出。定义
\begin{equation*}
 I_2(\Delta)
 \equiv\frac1{\ii}\mu_{\rm DR}^{2\epsilon}
 \int\frac{\dd^d r}{(2\pi)^d}
 \frac1{(r^2-\Delta+\ii0)^2}.
\end{equation*}
Wick 转动或使用标准 Gamma 函数积分给
\begin{align*}
 I_2(\Delta)
 &=\frac{\mu_{\rm DR}^{2\epsilon}}
 {(4\pi)^{d/2}}
 \Gamma\!\left(2-\frac d2\right)
 (\Delta-\ii0)^{d/2-2}\\
 &=\frac1{(4\pi)^2}
 \left[
 \frac1{\bar\epsilon}
 -\ln\frac{\Delta-\ii0}{\mu_{\rm DR}^2}
 +O(\epsilon)
 \right],
\end{align*}
其中
\begin{equation*}
 \frac1{\bar\epsilon}
 \equiv\frac1\epsilon-\gamma_E+\ln4\pi.
\end{equation*}
因此紫外极点与 $m_f$ 和 $q^2$ 无关；这些尺度只进入有限对数。代回并使用
\begin{equation*}
 \int_0^1\dd x\,x(1-x)=\frac16,
\end{equation*}
得到单个 Dirac 表示对正文所采用背景动能紫外系数的贡献
\begin{equation*}
 \Pi_{f,ab}^{\mu\nu}(q)\Big|_{\rm UV}
 =\delta_{ab}(q^2\eta^{\mu\nu}-q^\mu q^\nu)
 \frac{g_s^2}{(4\pi)^2}
 \left[-\frac43T(R_f)\right]
 \frac1{\bar\epsilon}.
\end{equation*}
对所有活跃 Dirac 表示求和即得
\begin{equation*}
 -\frac43\sum_fT(R_f),
\end{equation*}
也就是正文\eqnrefs{eq:qcd-quark-loop-uv-coefficient-dp10} 的一般物质屏蔽项；QCD 基本表示 $R_f=F$ 只在最后一步代入。
\end{derivation}
